\documentclass[11pt,a4paper]{article}
\usepackage[utf8]{inputenc}
\usepackage[T1]{fontenc}
\usepackage{amsfonts}
\usepackage[french]{babel}
\frenchbsetup{StandardLists=true}
\usepackage{enumitem}
\usepackage{amssymb}
\usepackage{amsmath}
\usepackage[left=2cm,right=2cm,top=2cm,bottom=2cm]{geometry}
\usepackage{multirow}
\usepackage[dvipsnames]{xcolor}
\usepackage{fouriernc}
\usepackage{graphicx}
\usepackage{setspace}
\singlespacing
\usepackage{tcolorbox}
\usepackage{multicol}
\usepackage{caption}
\usepackage{fancyhdr}
\pagestyle{fancy}
\renewcommand\headrulewidth{1pt}
\fancyhead[L]{Lycée Denis Diderot}
\fancyhead[R]{Classe de 3PM}
\setcounter{page}{1}\fancyfoot[C]{page \thepage}
\fancyfoot[R]{Année 2025-2026}
\fancyfoot[L]{Alexis RUIZ}
\usepackage{multirow,tabularx,booktabs}
\usepackage{pifont}
\usepackage{chemmacros}
\chemsetup{modules=all}
\setlength{\columnseprule}{.25mm}
\setlength{\parindent}{0pt}
\renewcommand{\arraystretch}{2}
\frenchbsetup{StandardLists=true}

\begin{document}
	
	\begin{center}
		\LARGE{CORRIGÉ -- Chapitre 4 -- Les ions}
		
		\vspace{3mm}
		
		\large{Durée : 50 minutes \hspace{1cm} \textbf{Barème total : /30 points}}
	\end{center}
	
	\normalsize
	
	\hrule
	
	
	\vspace{3mm}
	
	\hrule
	\vspace{3mm}
	
	\textbf{Exercice 1 : Connaissances de cours} \hfill \textbf{/6 points}
	
	\vspace{3mm}
	
	\textbf{1.} Qu'est-ce qu'un ion ? \hfill \textit{(1 pt)}
	
	\begin{tcolorbox}[colback=green!5!white,colframe=green!75!black]
		Un ion est un atome (ou un groupe d'atomes) qui a gagné ou perdu un ou plusieurs électrons. Il n'est donc plus électriquement neutre et porte une charge électrique.
	\end{tcolorbox}
	
	\vspace{3mm}
	
	\textbf{2.} Quelle est la différence entre un cation et un anion ? \hfill \textit{(2 pts)}
	
	\begin{tcolorbox}[colback=green!5!white,colframe=green!75!black]
		\begin{itemize}
			\item Un \textbf{cation} est un ion chargé positivement (+). Il a \textbf{perdu} des électrons. \textit{(1 pt)}
			\item Un \textbf{anion} est un ion chargé négativement (--). Il a \textbf{gagné} des électrons. \textit{(1 pt)}
		\end{itemize}
	\end{tcolorbox}
	
	\vspace{3mm}
	
	\textbf{3.} Comment un atome se transforme-t-il en ion ? \hfill \textit{(1 pt)}
	
	\begin{tcolorbox}[colback=green!5!white,colframe=green!75!black]
		Un atome se transforme en ion en gagnant ou en perdant un ou plusieurs électrons. Le nombre de protons dans le noyau ne change pas.
	\end{tcolorbox}
	
	\vspace{3mm}
	
	\textbf{4.} Citer deux ions que l'on peut identifier par des tests chimiques. \hfill \textit{(1 pt)}
	
	\begin{tcolorbox}[colback=green!5!white,colframe=green!75!black]
		Exemples : ions chlorure \ch{Cl-}, ions cuivre (II) \ch{Cu^{2+}}, ions fer (III) \ch{Fe^{3+}}, ions fer (II) \ch{Fe^{2+}}, ions aluminium \ch{Al^{3+}}... \textit{(0,5 pt par ion cité)}
	\end{tcolorbox}
	
	\vspace{3mm}
	
	\textbf{5.} Vrai ou Faux : Une solution ionique est toujours électriquement neutre. \hfill \textit{(1 pt)}
	
	\begin{tcolorbox}[colback=green!5!white,colframe=green!75!black]
		\textbf{Vrai} -- Une solution ionique est électriquement neutre car la somme des charges positives est égale à la somme des charges négatives.
	\end{tcolorbox}
	
	\vspace{5mm}
	
	\hrule
	\vspace{3mm}
	
	\textbf{Exercice 2 : Composition des ions} \hfill \textbf{/5 points}
	
	\vspace{3mm}
	
	L'atome d'aluminium possède 13 protons et 13 électrons. Il peut former l'ion aluminium \ch{Al^{3+}}.
	
	\vspace{3mm}
	
	\textbf{1.} L'ion aluminium est-il un cation ou un anion ? Justifier. \hfill \textit{(1 pt)}
	
	\begin{tcolorbox}[colback=green!5!white,colframe=green!75!black]
		L'ion aluminium \ch{Al^{3+}} est un \textbf{cation} car il porte une charge positive (+3). \textit{(0,5 pt pour la réponse + 0,5 pt pour la justification)}
	\end{tcolorbox}
	
	\vspace{2mm}
	
	\textbf{2.} Combien de protons possède l'ion aluminium \ch{Al^{3+}} ? \hfill \textit{(1 pt)}
	
	\begin{tcolorbox}[colback=green!5!white,colframe=green!75!black]
		L'ion aluminium possède \textbf{13 protons}. Le nombre de protons ne change pas lors de la formation d'un ion.
	\end{tcolorbox}
	
	\vspace{2mm}
	
	\textbf{3.} Combien d'électrons possède l'ion aluminium \ch{Al^{3+}} ? \hfill \textit{(1,5 pts)}
	
	\begin{tcolorbox}[colback=green!5!white,colframe=green!75!black]
		L'ion aluminium possède \textbf{10 électrons}.
		
		\textbf{Calcul :} 13 électrons (atome neutre) $-$ 3 électrons (perdus) $=$ 10 électrons. \textit{(0,5 pt pour le calcul + 1 pt pour la réponse)}
	\end{tcolorbox}
	
	\vspace{2mm}
	
	\textbf{4.} L'atome d'aluminium a-t-il gagné ou perdu des électrons pour devenir un ion ? Combien ? \hfill \textit{(1,5 pts)}
	
	\begin{tcolorbox}[colback=green!5!white,colframe=green!75!black]
		L'atome d'aluminium a \textbf{perdu 3 électrons} pour devenir l'ion \ch{Al^{3+}}. \textit{(1 pt pour "perdu" + 0,5 pt pour "3")}
	\end{tcolorbox}
	
	\vspace{5mm}
	
	
	\hrule
	\vspace{3mm}
	
	\textbf{Exercice 3 : Compléter le tableau} \hfill \textbf{/5 points}
	
	\vspace{3mm}
	
	\begin{center}
		\begin{tabular}{|l|c|c|c|c|}
			\hline
			\textbf{Particule} & \textbf{Nb de protons} & \textbf{Nb d'électrons} & \textbf{Atome ou ion ?} & \textbf{Formule} \\
			\hline
			Calcium & 20 & 20 & Atome & \ch{Ca} \\
			\hline
			Ion calcium & 20 & 18 & Ion & \ch{Ca^{2+}} \\
			\hline
			Oxygène & 8 & 8 & Atome & \ch{O} \\
			\hline
			Ion oxyde & 8 & 10 & Ion & \ch{O^{2-}} \\
			\hline
			Ion cuivre (II) & 29 & 27 & Ion & \ch{Cu^{2+}} \\
			\hline
		\end{tabular}
	\end{center}
	
	\begin{tcolorbox}[colback=blue!5!white,colframe=blue!75!black]
		\textbf{Barème :} 1 point par ligne correctement complétée (0,25 pt par case).
	\end{tcolorbox}
	
	\vspace{5mm}
	
	\hrule
	\vspace{3mm}
	
	\textbf{Exercice 4 : Tests de reconnaissance des ions} \hfill \textbf{/6 points}
	
	\vspace{3mm}
	
	\textbf{1.} Quel ion est présent dans la solution A (précipité rouille avec NaOH) ? \hfill \textit{(2 pts)}
	
	\begin{tcolorbox}[colback=green!5!white,colframe=green!75!black]
		La solution A contient des ions \textbf{fer (III)} \ch{Fe^{3+}}. \textit{(1 pt)}
		
		\textbf{Justification :} Le test avec l'hydroxyde de sodium produit un précipité rouille (brun-orangé), caractéristique des ions \ch{Fe^{3+}}. \textit{(1 pt)}
	\end{tcolorbox}
	
	\vspace{2mm}
	
	\textbf{2.} Quel ion est présent dans la solution B (précipité blanc qui noircit avec \ch{AgNO3}) ? \hfill \textit{(2 pts)}
	
	\begin{tcolorbox}[colback=green!5!white,colframe=green!75!black]
		La solution B contient des ions \textbf{chlorure} \ch{Cl-}. \textit{(1 pt)}
		
		\textbf{Justification :} Le test avec le nitrate d'argent produit un précipité blanc qui noircit à la lumière, caractéristique des ions chlorure \ch{Cl-}. \textit{(1 pt)}
	\end{tcolorbox}
	
	\vspace{2mm}
	
	\textbf{3.} Quel ion est présent dans la solution C (précipité bleu avec NaOH) ? \hfill \textit{(2 pts)}
	
	\begin{tcolorbox}[colback=green!5!white,colframe=green!75!black]
		La solution C contient des ions \textbf{cuivre (II)} \ch{Cu^{2+}}. \textit{(1 pt)}
		
		\textbf{Justification :} Le test avec l'hydroxyde de sodium produit un précipité bleu, caractéristique des ions cuivre (II). \textit{(1 pt)}
	\end{tcolorbox}
	
	\vspace{5mm}
	
	\hrule
	\vspace{3mm}
	
	\textbf{Exercice 5 : Identifier les bons tests} \hfill \textbf{/4 points}
	
	\vspace{3mm}
	
	Compléter le tableau des tests de reconnaissance :
	
	\vspace{3mm}
	
	\begin{center}
		\begin{tabular}{|l|c|c|}
			\hline
			\textbf{Ion à identifier} & \textbf{Réactif utilisé} & \textbf{Résultat positif} \\
			\hline
			Ion chlorure \ch{Cl-} & Nitrate d'argent \ch{AgNO3} & Précipité blanc qui noircit à la lumière \\
			\hline
			Ion cuivre (II) \ch{Cu^{2+}} & Hydroxyde de sodium \ch{NaOH} & Précipité bleu \\
			\hline
			Ion fer (II) \ch{Fe^{2+}} & Hydroxyde de sodium \ch{NaOH} & Précipité vert \\
			\hline
			Ion fer (III) \ch{Fe^{3+}} & Hydroxyde de sodium \ch{NaOH} & Précipité rouille (brun-orangé) \\
			\hline
		\end{tabular}
	\end{center}
	
	\begin{tcolorbox}[colback=blue!5!white,colframe=blue!75!black]
		\textbf{Barème :} 1 point par ligne correctement complétée (0,5 pt par case réactif/résultat).
	\end{tcolorbox}
	
	\vspace{5mm}
	
	\hrule
	\vspace{3mm}
	
	\textbf{Exercice 6 : Analyse de l'ion chlorure} \hfill \textbf{/4 points}
	
	\vspace{3mm}
	
	L'atome de chlore possède 17 protons et 17 électrons. Il peut former l'ion chlorure \ch{Cl-}.
	
	\vspace{3mm}
	
	\textbf{1.} Combien de protons et d'électrons possède l'ion chlorure \ch{Cl-} ? \hfill \textit{(2 pts)}
	
	\begin{tcolorbox}[colback=green!5!white,colframe=green!75!black]
		L'ion chlorure \ch{Cl-} possède :
		\begin{itemize}
			\item \textbf{17 protons} (le nombre de protons ne change jamais lors de la formation d'un ion) \textit{(1 pt)}
			\item \textbf{18 électrons} (17 + 1 électron gagné) \textit{(1 pt)}
		\end{itemize}
	\end{tcolorbox}
	
	\vspace{3mm}
	
	\textbf{2.} L'ion chlorure est-il un cation ou un anion ? Justifier. \hfill \textit{(2 pts)}
	
	\begin{tcolorbox}[colback=green!5!white,colframe=green!75!black]
		L'ion chlorure \ch{Cl-} est un \textbf{anion}. \textit{(1 pt)}
		
		\textbf{Justification :} Il porte une charge négative (--1) car il a gagné un électron. Les ions chargés négativement sont appelés anions. \textit{(1 pt)}
	\end{tcolorbox}
	
	\vspace{10mm}
	
	\begin{center}
		\LARGE{\textbf{Total : /30 points}}
	\end{center}
	
	\vspace{5mm}
	
	\begin{tcolorbox}[colback=red!5!white,colframe=red!75!black]
		\textbf{Conseils de correction :}
		\begin{itemize}
			\item Accepter les formulations équivalentes si elles sont scientifiquement correctes.
			\item Valoriser les justifications cohérentes même si la formulation n'est pas parfaite.
			\item Pour les tests : accepter \og soude \fg{} à la place de \og hydroxyde de sodium \fg{}.
			\item Pour l'exercice 3 : accepter les noms d'ions avec ou sans précision (ex. ion calcium ou \ch{Ca^{2+}}).
			\item Pour l'exercice 4 : accepter \og précipité brun \fg{} ou \og précipité orangé \fg{} pour le fer (III).
		\end{itemize}
	\end{tcolorbox}
	
\end{document}